\documentclass[12pt, oneside]{book}

\usepackage{graphicx}
\usepackage{amsmath}
\usepackage{amsfonts}
\usepackage{amssymb}
\usepackage{hyperref}
\usepackage{caption}

\usepackage[utf8]{inputenc}
\usepackage[T2A,T1]{fontenc}
\usepackage{ragged2e}

\usepackage{setspace}
\setstretch{1.2}
\setlength{\parindent}{1.25cm}
\setlength{\emergencystretch}{3em}
\everymath{\displaystyle}
\usepackage{tocloft}
\cftsetindents{section}{1em}{2em}
\renewcommand\cfttoctitlefont{\hfill\Large\bfseries}
\renewcommand\cftaftertoctitle{\hfill\mbox{}}
\setlength{\cftbeforesecskip}{0.1cm}
\setlength{\cftbeforesubsecskip}{0.1cm}
\setlength{\cftbeforetoctitleskip}{0cm}
\setlength{\cftaftertoctitleskip}{0.4cm}
\setcounter{tocdepth}{2}

\usepackage[english,russian]{babel}
\usepackage[a4paper, left=3cm, top=2cm, right=1cm, bottom=2cm]{geometry}

\makeatletter
\renewcommand{\@makeschapterhead}[1]{%
        \vspace*{0\p@}
        {\parindent \z@ \center
                \normalfont
                \LARGE \bfseries #1\par\nobreak
                \addcontentsline{toc}{chapter}{#1}
                \vskip 10\p@
}}
\makeatother

\makeatletter
\renewcommand{\chapter}{%
        \thispagestyle{plain}
        \global\@topnum\z@
        \@afterindentfalse
\secdef\@chapter\@schapter}

\renewcommand{\@makechapterhead}[1]{%
        \vspace*{0\p@}
        {\parindent \z@ \raggedright
                \normalfont
                \LARGE \bfseries \thechapter\quad #1\par\nobreak
                \vskip 20\p@
}}
\makeatother

\makeatletter
\newcommand{\custompagestyle}{%
        \begingroup
        \renewcommand{\ps@plain}{%
                \renewcommand{\@oddfoot}{\hfil\thepage\hfil}
                \renewcommand{\@evenfoot}{\hfil\thepage\hfil}
                \renewcommand{\@oddhead}{}
                \renewcommand{\@evenhead}{}
        }%
        \pagestyle{plain}
}
\makeatother

\newenvironment{customquote}
{\list{}{\leftmargin=0.5cm
                \rightmargin=0.5cm
                \topsep=0pt
                \partopsep=0pt
                \parsep=2pt
                \itemsep=0pt
        }%
        \normalfont\relax
\item\relax}
{\endlist}

\begin{document}
\thispagestyle{empty}
\begin{center}
        \textbf{\large САНКТ-ПЕТЕРБУРГСКИЙ ГОСУДАРСТВЕННЫЙ УНИВЕРСИТЕТ} \\[2cm]

        {\large ОТЧЕТ О НАУЧНО-ИССЛЕДОВАТЕЛЬСКОЙ РАБОТЕ} \\[2cm]

        {\Large Шаго Павел Евгеньевич} \\[0.3cm]

        {22.Б07-ПУ, 01.03.02 Прикладная математика и информатика} \\[1cm]

        \textbf{\Large Планировщики 4} \\[3cm]

        {\large Научный руководитель} \\[0.3cm]
        \large{к.ф.-м.н., доцент Корхов В. В.} \\[10cm]
        {Санкт-Петербург}\\ \today

\end{center}

\newpage
\custompagestyle
\setcounter{page}{2}
\renewcommand{\contentsname}{Содержание}
\tableofcontents

\newpage
\chapter*{Введение}
\begin{customquote}
        \quad Earliest Eligible Virtual Deadline First
        (EEVDF)~\cite{stoica} --- алгоритм
        пропорцио\-наль\-но-доле\-вого планирования,
        принятый в качестве основного планировщика ядра
        Linux начиная с версии 6.6~\cite{eevdf_intro}
        взамен Completely Fair Scheduler
        (CFS)~\cite{wong}. Фреймворк
        sched\_ext~\cite{schedext} (с версии Linux 6.12)
        позволяет загружать пользовательские планировщики
        как eBPF-программы без пересборки ядра, что
        открывает возможность прототипирования
        алгоритмов~\cite{srand}, однако вносит
        накладные расходы и ограничивает доступ к ряду
        подсистем ядра.

        \quad Целью данной работы является исследование
        возможности реализации алгоритма EEVDF средствами
        фреймворка sched\_ext. Для этого приводится
        математическая модель алгоритма, строится явное
        отображение каждого шага модели на конкретный
        обработчик sched\_ext и проводится экспериментальная
        оценка с использованием специализированных тестов производительности
        (hackbench, sysbench, schbench) на трёх уровнях
        нагрузки. Результаты позволяют оценить накладные
        расходы фреймворка sched\_ext и компромиссы
        производительности при переносе алгоритма
        из ядра во фреймворк.
\end{customquote}

\newpage
\chapter{Математическая модель EEVDF}

\section{Пропорционально-долевое планирование}
\begin{customquote}
        \quad Earliest Eligible Virtual Deadline First
        (EEVDF)~\cite{stoica} --- модель
        пропорцио\-наль\-но-доле\-вого планирования, предложенная
        Ion Stoica и Hussein Abdel-Wahab.

        \quad Рассматривается один разделяемый ресурс
        (процессор), время на котором выделяется квантами
        длительности не более $q$. Множество клиентов (задач),
        активных в момент $t$, обозначается $\mathcal{A}(t)$,
        каждому клиенту $i$ назначен положительный вес $w_i$.
        Идеальная мгновенная доля ресурса для клиента $i$:
        \begin{equation}
                f_i(t) = \frac{w_i}{\sum_{j \in \mathcal{A}(t)} w_j}.
        \end{equation}
        В идеализированной непрерывной модели при $q \to 0$
        обслуживание, положенное клиенту $i$ на интервале
        $[t_0, t_1]$, есть интеграл от $f_i$:
        \begin{equation}
                S_i(t_0, t_1) = \int_{t_0}^{t_1} f_i(\tau)\, d\tau.
        \end{equation}
        Разность между положенным и фактически полученным
        обслуживанием задаёт отставание (lag):
        \begin{equation}
                \mathrm{lag}_i(t) = S_i(t_0^i, t) - s_i(t_0^i, t),
        \end{equation}
        где $t_0^i$ --- момент входа клиента $i$ в систему,
        $s_i$ --- фактически полученное обслуживание.
        Положительный lag соответствует недообслуженному
        клиенту, отрицательный --- переобслуженному.

        \quad Виртуальное время системы $V(t)$:
        \begin{equation}
                V(t) = \int_0^t \frac{d\tau}{\sum_{j \in
                \mathcal{A}(\tau)} w_j}.
        \end{equation}
        За одну единицу виртуального времени каждый активный
        клиент получает ровно $w_i$ единиц реального
        обслуживания:
        \begin{equation}
                S_i(t_1, t_2) = w_i \cdot \bigl(V(t_2) - V(t_1)\bigr).
        \end{equation}
\end{customquote}

\section{Виртуальные дедлайны и правило выбора}
\begin{customquote}
        \quad Каждый запрос $k$ клиента $i$ характеризуется
        длиной $r^{(k)}$, виртуальным временем доступности
        $ve^{(k)}$ и виртуальным дедлайном $vd^{(k)}$.
        Время доступности:
        \begin{equation}
                ve^{(k)} = V(t_0^i) + \frac{s_i(t_0^i, t)}{w_i},
        \end{equation}
        виртуальный дедлайн:
        \begin{equation}
                vd^{(k)} = ve^{(k)} + \frac{r^{(k)}}{w_i}.
        \end{equation}
        Запрос считается доступным (eligible) в момент $t$,
        если $ve \leq V(t)$. Среди доступных запросов
        выбирается тот, чей $vd$ минимален.

        \quad При полном использовании кванта
        $ve^{(k+1)} = vd^{(k)}$; при частичном использовании
        $u^{(k)} \leq r^{(k)}$ обновление принимает вид
        $ve^{(k+1)} = ve^{(k)} + u^{(k)} / w_i$.

        \textbf{Лемма 1} (доступность недообслуженного
        клиента)~\cite{stoica}.
        \textit{Если $\mathrm{lag}_i(t) \geq 0$ для активного
        клиента $i$, то $i$ имеет доступный запрос в $t$.}

        \textbf{Лемма 2}~\cite{stoica}.
        \textit{В любой момент $t$ сумма отставаний равна нулю:}
        \begin{equation}
                \sum_{i \in \mathcal{A}(t)} \mathrm{lag}_i(t) = 0.
        \end{equation}
        Из лемм 1 и 2 следует, что EEVDF является
        сохраняющим работу (work-conserving): ресурс не
        простаивает при наличии активных клиентов.
\end{customquote}

\section{Границы отставания}
\begin{customquote}
        \textbf{Теорема 1} (границы
        отставания)~\cite{stoica}.
        \textit{В стационарной системе (множество активных
        клиентов не меняется) для любого активного
        клиента $i$:}
        \begin{equation}
                -r_{\max} < \mathrm{lag}_i(t) < \max(r_{\max}, q),
        \end{equation}
        \textit{где $r_{\max}$ --- максимальная длина запроса
        клиента $i$. Обе границы асимптотически точны.}

        \textbf{Следствие.}
        \textit{Если $r^{(k)} \leq q$ для всех запросов, то
        $-q < \mathrm{lag}_i(t) < q$.}

        \textbf{Лемма} (оптимальность)~\cite{stoica}.
        \textit{Ни один пропорцио\-наль\-но-доле\-вой алгоритм
        с квантами размера $q$ не может дать границу лучше
        $(-q, q)$, и EEVDF достигает этой границы.}

        \quad Фильтр доступности ($ve \leq V(t)$) отличает
        EEVDF от алгоритмов справедливой очередизации
        (Fair Queueing)~\cite{leung}, ограничивая отставание
        до $O(1)$ вместо $O(n)$~\cite{stoica}.
\end{customquote}

\chapter{Отображение модели на sched\_ext}

\begin{customquote}
        \quad Описание фреймворка sched\_ext~\cite{schedext}
        и его обратных вызовов (ops) приведено в
        предыдущих работах~\cite{nir2,nir3}. Ниже
        описываются переменные состояния и отображение
        обработчиков, специфичные для реализации EEVDF.
\end{customquote}

\section{Переменные состояния}
\begin{customquote}
        \quad Глобальное состояние хранится в eBPF-карте с
        единственным элементом: \texttt{vtime\_now} ---
        системное виртуальное время $V(t)$, масштабированное
        коэффициентом $S = 1000$ для целочисленной точности,
        и \texttt{total\_weight} --- суммарный вес $\sum w_j$.

        \quad Состояние задачи хранится в карте
        \texttt{BPF\_MAP\_TYPE\_TASK\_STORAGE} с обращением
        за $O(1)$: виртуальное время доступности \texttt{ve}
        ($ve^{(k)}$, уравнение~6), виртуальный дедлайн
        \texttt{vd} ($vd^{(k)}$, уравнение~7).
        Единственная очередь \texttt{SHARED\_DSQ} упорядочена
        по $vd$: вставка через
        \texttt{scx\_bpf\_dsq\_insert\_vtime} реализует
        правило выбора EEVDF --- среди доступных запросов
        извлекается тот, чей $vd$ минимален.
\end{customquote}

\section{Отображение обработчиков}
\begin{customquote}
        \quad В таблице~\ref{tab:eevdf_ops} приведено
        соответствие между обработчиками sched\_ext и
        шагами модели EEVDF.

        \begin{table}[!ht]
                \centering
                \caption{Отображение обработчиков
                        sched\_ext на модель EEVDF}
                \label{tab:eevdf_ops}
                \vspace{0.3em}
                \begin{tabular}{|l|p{9.5cm}|}
                        \hline
                        \textbf{Обработчик} &
                                \textbf{Реализуемый шаг} \\
                        \hline
                        \texttt{select\_cpu} &
                                Быстрый путь: если
                                $ve \leq V(t)$ и есть свободное
                                ядро, задача помещается в
                                \texttt{LOCAL\_DSQ} \\
                        \texttt{enqueue} &
                                Ограничитель отставания
                                (теорема~1), вычисление
                                $vd = ve + q \cdot S / w_i$,
                                вставка в
                                \texttt{SHARED\_DSQ} по $vd$ \\
                        \texttt{dispatch} &
                                Извлечение головы DSQ ---
                                задачи с наименьшим $vd$ \\
                        \texttt{stopping} &
                                $ve \mathrel{+}= u / w_i$,
                                $V(t) \mathrel{+}=
                                u \cdot S / \sum w_j$ \\
                        \texttt{enable} &
                                Вход: $ve := V(t)$,
                                $\sum w_j \mathrel{+}= w_i$ \\
                        \texttt{disable} &
                                Выход: $V(t^+) = V(t) +
                                \mathrm{lag}/\sum w_{j \neq i}$ \\
                        \texttt{set\_weight} &
                                Изменение веса:
                                выход--вход с новым весом \\
                        \hline
                \end{tabular}
        \end{table}

        \quad \textbf{Атомарное продвижение $V(t)$.}
        На многоядерной системе несколько ядер могут
        одновременно вызвать \texttt{stopping}. Продвижение
        $V$ выполняется через атомарную операцию
        \texttt{\_\_sync\_fetch\_and\_add}:
        \begin{equation}
                \Delta V = \frac{u \cdot S}{\sum_{j} w_j},
        \end{equation}
        где $u$ --- фактически использованная часть кванта.

        \quad \textbf{Двусторонний ограничитель отставания.}
        При постановке задачи в очередь значение $ve$
        корректируется до $ve'$:
        \[
                \max(V(t) - q_v,\; ve) \leq ve'
                \leq V(t) + q_v,
        \]
        где $q_v = q \cdot S / w_i$. Нижняя граница
        предотвращает монополизацию ресурса задачей с
        накопленным отставанием; верхняя не позволяет задаче
        уйти слишком далеко в будущее. Оба ограничения
        следуют из границ теоремы~1.

        \quad Единственное отклонение от теоретической модели ---
        замена непрерывного продвижения $V(t)$ дискретными
        приращениями на границах квантов. Это не нарушает
        границ теоремы~1, поскольку $q$ конечно.
\end{customquote}

\chapter{Экспериментальная оценка}

\section{Тестовое оборудование и методика}
\begin{customquote}
        \quad Экспериментальная оценка проведена на платформе
        с процессором Intel Core Ultra~7 265K
        (Arrow Lake-S)~\cite{intelhybrid}: 8 P-ядер
        Lion Cove (до 5,5~ГГц) и 12 E-ядер Skymont
        (до 4,6~ГГц), 20 ядер без Hyper-Threading,
        32 ГБ DDR5. Операционная
        система --- Linux 7.0 с поддержкой sched\_ext.

        \quad Сравниваются штатный EEVDF ядра Linux
        (эталон) и его sched\_ext реализация.
        Набор тестов:
        \begin{itemize}
                \item \textbf{hackbench} --- массовое порождение
                        задач с интенсивным IPC через каналы.
                        Метрика: время выполнения (секунды).
                \item \textbf{sysbench} --- OLTP-нагрузка
                        (чтение PostgreSQL, 4 таблицы по
                        100\,000 строк, 10\,с).
                        Метрика: TPS.
                \item \textbf{schbench} --- чувствительная к задержкам
                        веб-нагрузка. Метрики: задержка
                        пробуждения (p99, p99.9), RPS.
        \end{itemize}

        \quad Параметры нагрузки приведены в
        таблице~\ref{tab:params}.

        \begin{table}[!ht]
                \centering
                \caption{Параметры тестов по уровням нагрузки}
                \label{tab:params}
                \vspace{0.3em}
                \begin{tabular}{|l|c|c|c|}
                        \hline
                        Тест & Лёгкая & Умеренная & Стрессовая \\
                        \hline
                        hackbench & \texttt{-l 10000}
                                  & \texttt{-g 4 -l 120000}
                                  & \texttt{-g 10 -l 50000} \\
                        sysbench (потоки) & 1 & 4 & 20 \\
                        schbench (\texttt{-m}/\texttt{-t})
                                  & 1\,/\,5 & 2\,/\,20
                                  & 4\,/\,40 \\
                        \hline
                \end{tabular}
        \end{table}

        \quad Каждый тест запускается при трёх уровнях
        нагрузки по три прогона со случайным порядком
        планировщиков. Приводятся средние значения и
        95\%-доверительные интервалы по критерию Стьюдента.
\end{customquote}

\section{Результаты}
\begin{customquote}
        \quad \textbf{hackbench}
        (таблица~\ref{tab:hackbench}).

        \begin{table}[!ht]
                \centering
                \caption{hackbench: время выполнения (секунды)}
                \label{tab:hackbench}
                \vspace{0.3em}
                \begin{tabular}{|l|c|c|c|c|c|c|}
                        \hline
                        Планиров- & \multicolumn{2}{c|}{Лёгкая}
                                  & \multicolumn{2}{c|}{Умеренная}
                                  & \multicolumn{2}{c|}{Стрессовая} \\
                        \cline{2-7}
                        щик & Ср. & 95\%-ДИ & Ср. & 95\%-ДИ & Ср. & 95\%-ДИ \\
                        \hline
                        sched\_ext EEVDF & \textbf{1,838} & {\scriptsize [1,752;\:1,923]}
                                   & \textbf{8,716} & {\scriptsize [8,640;\:8,793]}
                                   & \textbf{9,176} & {\scriptsize [8,998;\:9,354]} \\
                        Linux EEVDF & 1,900 & {\scriptsize [1,870;\:1,931]}
                                   & 9,655 & {\scriptsize [9,628;\:9,683]}
                                   & 9,344 & {\scriptsize [9,189;\:9,498]} \\
                        \hline
                \end{tabular}
        \end{table}

        \quad Средние значения sched\_ext ниже штатного
        на 1,8--9,7\%, но доверительные интервалы
        пересекаются при лёгкой и стрессовой нагрузке ---
        статистически значимое различие только при
        умеренной (9,7\%). Реализация через sched\_ext
        не вносит измеримого ухудшения для данного типа
        нагрузки.

        \quad \textbf{sysbench}
        (таблица~\ref{tab:sysbench}).

        \begin{table}[!ht]
                \centering
                \caption{sysbench: транзакции в секунду (TPS)}
                \label{tab:sysbench}
                \vspace{0.3em}
                \begin{tabular}{|l|c|c|c|c|c|c|}
                        \hline
                        Планиров- & \multicolumn{2}{c|}{Лёгкая}
                                  & \multicolumn{2}{c|}{Умеренная}
                                  & \multicolumn{2}{c|}{Стрессовая} \\
                        \cline{2-7}
                        щик & Ср. & 95\%-ДИ & Ср. & 95\%-ДИ & Ср. & 95\%-ДИ \\
                        \hline
                        Linux EEVDF & \textbf{4246} & {\scriptsize [4201;\:4292]}
                                   & \textbf{16\,566} & {\scriptsize [16\,487;\:16\,645]}
                                   & \textbf{67\,373} & {\scriptsize [65\,042;\:69\,705]} \\
                        sched\_ext EEVDF & 3793 & {\scriptsize [3295;\:4292]}
                                   & 12\,517 & {\scriptsize [12\,279;\:12\,756]}
                                   & 41\,803 & {\scriptsize [39\,324;\:44\,282]} \\
                        \hline
                \end{tabular}
        \end{table}

        \quad На sysbench штатный EEVDF опережает
        sched\_ext на 10,7--38,0\% по TPS (при лёгкой
        нагрузке ДИ пересекаются). Снижение обусловлено
        накладными расходами eBPF-вызовов при высоком
        темпе переключений контекста и отсутствием
        интеграции с подсистемой wake\_affine.

        \quad \textbf{schbench}
        (таблица~\ref{tab:schbench}).

        \begin{table}[!ht]
                \centering\small
                \caption{schbench: задержка
                        пробуждения (мкс) и RPS}
                \label{tab:schbench}
                \vspace{0.3em}
                \begin{tabular}{|l|c|c|c|c|c|c|}
                        \hline
                        Планиров- & \multicolumn{2}{c|}{p99}
                                  & \multicolumn{2}{c|}{p99.9}
                                  & \multicolumn{2}{c|}{Ср. RPS} \\
                        \cline{2-7}
                        щик & Ср. & 95\%-ДИ & Ср. & 95\%-ДИ
                            & Ср. & 95\%-ДИ \\
                        \hline
                        \multicolumn{7}{|c|}{Лёгкая нагрузка} \\
                        \hline
                        sched\_ext EEVDF & \textbf{5} & ---
                                  & \textbf{59} & {\scriptsize [54;\:64]}
                                  & \textbf{1668} & {\scriptsize [1522;\:1814]} \\
                        Linux EEVDF & 152 & {\scriptsize [141;\:163]}
                                & 185 & {\scriptsize [182;\:189]}
                                & 591 & {\scriptsize [589;\:593]} \\
                        \hline
                        \multicolumn{7}{|c|}{Умеренная нагрузка} \\
                        \hline
                        sched\_ext EEVDF & \textbf{2375} & {\scriptsize [2277;\:2473]}
                                  & \textbf{2815} & {\scriptsize [2713;\:2917]}
                                  & \textbf{8585} & {\scriptsize [8582;\:8588]} \\
                        Linux EEVDF & 3300 & {\scriptsize [3266;\:3334]}
                                & 4733 & {\scriptsize [4281;\:5185]}
                                & 5700 & {\scriptsize [5673;\:5726]} \\
                        \hline
                        \multicolumn{7}{|c|}{Стрессовая нагрузка} \\
                        \hline
                        sched\_ext EEVDF & 16\,112 & {\scriptsize [16\,033;\:16\,191]}
                                  & \textbf{16\,672} & ---
                                  & \textbf{8639} & {\scriptsize [8628;\:8650]} \\
                        Linux EEVDF & \textbf{11\,131} & {\scriptsize [10\,534;\:11\,727]}
                                & 17\,163 & {\scriptsize [15\,266;\:19\,059]}
                                & 5645 & {\scriptsize [5434;\:5856]} \\
                        \hline
                \end{tabular}
        \end{table}

        \quad При лёгкой и умеренной нагрузке sched\_ext
        EEVDF превосходит штатный по всем метрикам
        благодаря агрессивному быстрому пути
        \texttt{select\_cpu}: при наличии свободного ядра
        задача помещается в \texttt{LOCAL\_DSQ}, минуя
        общую очередь. При стрессовой нагрузке sched\_ext
        проигрывает по p99, но выигрывает
        по RPS ($8639$ vs $5645$,
        разница 53\%) за счёт более равномерной загрузки
        ядер через глобальную очередь DSQ.

\end{customquote}


\newpage
\chapter*{Заключение}
\begin{customquote}
        \quad В работе приведена математическая модель
        EEVDF, построено её отображение на обработчики
        sched\_ext (таблица~\ref{tab:eevdf_ops}) и
        проведено экспериментальное сравнение со штатным
        EEVDF ядра Linux:
        \begin{itemize}
                \item \textbf{hackbench} (массовый IPC):
                        sched\_ext статистически значимо
                        быстрее при умеренной нагрузке
                        (на 9,7\%), при лёгкой и стрессовой
                        различия не значимы --- накладные
                        расходы фреймворка пренебрежимо малы
                        при данном типе нагрузки;
                \item \textbf{schbench} (веб-нагрузка):
                        sched\_ext превосходит
                        штатный EEVDF по RPS на всех уровнях
                        нагрузки (в 2,8 раза при лёгкой,
                        на 50,6\% при умеренной, на 53\%
                        при стрессовой);
                \item \textbf{sysbench} (OLTP): sched\_ext
                        уступает на 10--38\% по TPS из-за
                        стоимости eBPF-вызовов при высоком
                        темпе переключений и отсутствия
                        интеграции с wake\_affine.
        \end{itemize}

        \quad Результаты показывают, что sched\_ext
        пригоден для прототипирования алгоритмов
        планирования: реализация корректно воспроизводит
        поведение EEVDF на hackbench и schbench, но
        на нагрузках с интенсивным переключением контекста
        (sysbench) накладные расходы фреймворка могут быть
        существенными. Это ограничение носит архитектурный
        характер и не зависит от конкретного алгоритма,
        что следует учитывать при выборе sched\_ext
        в качестве платформы для исследований.
\end{customquote}

\newpage
\renewcommand{\bibname}{\large СПИСОК ИСПОЛЬЗОВАННЫХ ИСТОЧНИКОВ}
\begin{thebibliography}{99}
        \vspace*{0.3cm}
        \bibitem{stoica}
        Ion Stoica, Hussein Abdel-Wahab Earliest Eligible Virtual Deadline First
        : A Flexible and Accurate Mechanism for Proportional Share
        Resource Allocation
        Old Dominion University 1996
        \bibitem{schedext}
        Jonathan Corbet The BPF extensible scheduler class
        [Электронный ресурс].
        -- Режим доступа: \url{https://lwn.net/Articles/916291/} --
        Дата обращения 17.02.2026
        \bibitem{eevdf_intro}
        An EEVDF CPU scheduler for Linux [Электронный ресурс].
        -- Режим доступа: \url{https://lwn.net/Articles/925371} --
        Дата обращения 04.03.2026
        \bibitem{wong}
        Chee Siang Wong, Ian K. T. Tan, Rosalina D. Kumari,
        Joy W. Lam, Wee Fun
        Fairness and Interactive Performance of O(1) and CFS
        Linux Kernel Schedulers
        International Symposium on Information Technology 2008
        \bibitem{intelhybrid}
        Efraim Rotem et al. Alder Lake Architecture
        2021 IEEE Hot Chips 33 Symposium (HCS) Palo Alto CA
        USA 2021 pp. 1--23
        \bibitem{srand}
        Qianhui Tang, Xiao Gao, Guodong Li, Jian Lin
        Optimizing Linux Scheduling Based on Global Runqueue
        with SCX
        IEEE International Conference SMC 2024
        \bibitem{leung}
        Joseph Y-T. Leung\ Handbook of Scheduling: Algorithms, Models, and
        Performance Analysis CRC Press 2004
        \bibitem{nir2}
        Павел Е. Шаго Отчет о научно-исследовательской работе:
        Планировщики 2 б. и.
        2025
        \bibitem{nir3}
        Павел Е. Шаго Отчет о научно-исследовательской работе:
        Планировщики 3 б. и.
        2025
\end{thebibliography}

\end{document}
